\section{Introduction}
\label{sec:introduction}

The deployment of autonomous multi-agent systems into high-stakes operational environments — agricultural irrigation, financial trading, medical triage, infrastructure management — has exposed a fundamental gap in the architecture of modern AI governance. These systems coordinate complex actions across distributed nodes without any guaranteed mechanism to escalate accountability when conditions exceed their design envelope. The result is a class of failures we characterize as \textit{accountability orphaning}: the systematic disconnection of consequential actions from any human actor capable of authorizing, halting, or correcting them.

Accountability orphaning is not a hypothetical edge case. It is the structural consequence of deploying autonomous agents in environments where the combinatorial space of possible failure modes exceeds what any pre-programmed error handler can anticipate. A soil moisture sensor that fails in a way that systematically misreports dryness will cause over-irrigation with no alarm triggered. An autonomous vehicle that encounters a scenario outside its training distribution will take potentially dangerous actions with no human aware in time to intervene. A financial trading agent operating at millisecond latency will execute trades that no human reviewer can assess before consequences materialize. In each case, the system's failure mode is silent, the system's knowledge of its own failure is absent, and the human accountability chain is structurally broken.

Conventional approaches to this problem rely on static error-handling code: \texttt{if sensor X fails, do Y}. This approach fails at scale for two reasons. First, the combinatorial space of multi-point failure scenarios in distributed sensor-actuator networks exceeds what any engineering team can anticipate and code against explicitly. Static error handling is a lossy compression of the failure space — it handles the scenarios that were anticipated, and silently mishandles everything else. Second, static error handlers are themselves implemented in software and therefore subject to the same failure modes as the system they govern: they can be bypassed, misconfigured, or compromised by the same agentic processes they are meant to constrain.

A more fundamental problem underlies both failure modes. Current autonomous systems treat accountability as a \textit{policy property} — something enforced by configuration, access control lists, or operational procedures — rather than as an \textit{architectural property} with the same enforceability guarantees as physical constraint. When accountability is a policy, it can be overridden by a sufficiently authorized actor. When accountability is an architecture, it survives compromise, delegation, and network disconnection. The distinction is not subtle: a policy can be changed by the system it governs; an architecture cannot.

We propose the Bailout Protocol as a first-principles solution to this architectural gap. The Bailout Protocol is a mandatory vertical-escalation mechanism that guarantees a human resolver receives every high-stakes exception before any machine actor can act on it. It is defined not as a software configuration but as a structural property of the system's execution graph: once a bailout trigger condition fires, all machine actors are architecturally suspended and control is transferred directly to a human accountability anchor through the \fact layer, bypassing the \purpose and \bounds layers entirely.

The core insight driving this design is that accountability without an architectural fallback is not accountability — it is aspiration. Any system that claims human accountability but implements it as a configurable permission rather than a structurally enforced interrupt is vulnerable to the same failure modes it claims to prevent. The Bailout Protocol closes this gap by making human accountability a fallback that operates below the level at which machine actors can intervene.

We note that the Bailout Protocol operates orthogonally to the Training Wheels Protocol described in our prior work \cite{bx3-framework}. The Training Wheels Protocol governs the system's general mode of operation based on aggregate trust scores — HITL Active, Shadow Mode, Full Autonomy, or Lockdown. The Bailout Protocol is event-driven and governs the escalation path within any training wheels mode. Critically, a system running in Full Autonomy (Training Wheels Mode 0) remains subject to the Bailout Protocol: if the Truth Gate fails or a high-risk trigger condition fires, the bailout path activates and suspends all machine actors, including those operating under full autonomy. The Training Wheels Protocol's Lockdown mode is itself a terminal form of bailout, but the Bailout Protocol generalizes this concept to all escalation scenarios regardless of trust-based mode.

The contribution of this paper is a complete formal specification of the Bailout Protocol, including: the trigger condition taxonomy (risk threshold, Truth Gate failure, and autonomy boundary); a six-state finite state machine with formal transition predicates; an escalation algorithm with independent timeout handling; an architectural bypass mechanism that renders all machine actors moot once bailout is initiated; formal proofs of bypass soundness, non-suppression, and deterministic resolution; and an analysis of the protocol's relationship to the Training Wheels Protocol and the broader \bxthree Framework.

The remainder of this paper is organized as follows. Section~\ref{sec:motivation} provides motivating examples drawn from autonomous agricultural systems, financial trading, and autonomous vehicles that illustrate the failure modes the Bailout Protocol addresses. Section~\ref{sec:related} reviews related work in human-in-the-loop AI, exception handling in autonomous systems, and AI governance frameworks. Section~\ref{sec:protocol} presents the formal specification of the protocol: the state machine, transition predicates, trigger conditions, escalation algorithm, and bypass mechanism. Section~\ref{sec:properties} proves the three formal properties of the protocol. Section~\ref{sec:timeouts} analyzes timeout handling and clock independence. Section~\ref{sec:trainingwheels} clarifies the relationship between the Bailout Protocol and the Training Wheels Protocol. Section~\ref{sec:conclusion} concludes and discusses implications for the design of next-generation autonomous systems.

\medskip
\noindent\textit{Note: The first-person plural ``we'' is used in the conventional academic sense, consistent with single-author scholarly writing.}

\bibliographystyle{plain}
\bibliography{bx3framework}